% !TeX root = ../复分析笔记.tex
% 由主文档通过 \input 载入；请编译主文档。
\setcounter{chapter}{1} % 保留原书章号。
\chapter{复数与复平面的几何}
\label{chap:ca-plane}

\section{复数的运算、共轭与模}

\begin{definition}[复数与共轭]
形如 $z=x+iy$ 的数称为\emph{复数（complex number）}，其中 $i^2=-1$。$x$ 为实部，$y$ 为虚部，分别记作 $\operatorname{Re}z$、$\operatorname{Im}z$；虚部是实数 $y$，不是 $iy$。复数 $\overline z=x-iy$ 称为 $z$ 的\emph{共轭（conjugate）}。
\end{definition}

两个复数相等，当且仅当实部、虚部分别相等。加法按分量进行，乘法按分配律展开后使用 $i^2=-1$：
\[
 (x+iy)(u+iv)=(xu-yv)+i(xv+yu).
\]
因此，复数可以作为平面向量相加，但乘法包含了普通二维向量没有的额外结构。除法可借助共轭把分母化成实数；当 $u+iv\ne0$ 时，
\[
 \frac{x+iy}{u+iv}
 =\frac{(x+iy)(u-iv)}{u^2+v^2}
 =\frac{xu+yv}{u^2+v^2}+i\frac{yu-xv}{u^2+v^2}.
\]

\begin{definition}[模与距离]
$z=x+iy$ 的\emph{模（modulus）}定义为 $|z|=\sqrt{x^2+y^2}$。复平面中两点 $z,w$ 的距离为 $|z-w|$。特别地，$z\overline z=|z|^2$，故非零复数的倒数为 $z^{-1}=\overline z/|z|^2$。
\end{definition}

共轭与四则运算相容，例如 $\overline{zw}=\overline z\,\overline w$；模满足 $|zw|=|z||w|$。后一个等式说明模把复数乘法转化为非负实数乘法。另一方面，模不能与加法直接交换，正确的关系是下面的不等式。

\begin{proposition}[三角不等式]
对任意 $z,w\in\mathbb C$，
\begin{equation}
 |z+w|\le |z|+|w|,\qquad
 \bigl||z|-|w|\bigr|\le |z-w|.
 \label{eq:ca-triangle}
\end{equation}
\end{proposition}
\begin{proof}
由 $\operatorname{Re}(z\overline w)\le|z\overline w|=|z||w|$，有
\[
 |z+w|^2=|z|^2+2\operatorname{Re}(z\overline w)+|w|^2
 \le(|z|+|w|)^2.
\]
两边开平方得到第一个不等式。将它应用于 $z=(z-w)+w$，得到 $|z|-|w|\le|z-w|$；交换 $z,w$ 后再合并即得第二个不等式。
\end{proof}

几何上，第一个不等式就是三角形的两边之和不小于第三边；当两个非零向量 $z,w$ 同向时取等号。估计复杂表达式时，先用三角不等式把各项拆开，再使用模的乘法性质，通常比先分离实部与虚部更简洁。

\begin{example}[从复数方程读出几何轨迹]
求满足 $|z-1|=|z+1|$ 的点集，以及满足 $|z-1|<|z+1|$ 的点集。
\end{example}
\begin{solve}
令 $z=x+iy$。平方后分别比较 $(x-1)^2+y^2$ 与 $(x+1)^2+y^2$，等式给出 $x=0$，严格不等式给出 $x>0$。前者是线段 $[-1,1]$ 的垂直平分线，后者是到 $1$ 比到 $-1$ 更近的右半平面。距离表达式同时提供了代数算法和几何解释。
\end{solve}

复数本身不存在与域运算相容的全序：在有序域中非零元素的平方为正，但 $i^2=-1$ 会与 $1>0$ 矛盾；比较大小时应明确比较的是实数或复数的模。

\section{极形式、辐角与开方}

非零复数 $z$ 可以写成 $z=r(\cos\theta+i\sin\theta)$，其中 $r=|z|>0$。此时 $\theta$ 称为 $z$ 的\emph{辐角（argument）}；同一复数对应的全部辐角为 $\theta+2k\pi$，$k\in\mathbb Z$。记 $\arg z$ 为这组值，记 $\operatorname{Arg}z\in(-\pi,\pi]$ 为辐角主值。零没有辐角。主值区间只是约定，它在负实轴处存在跳跃，并不是整个去心平面上的连续函数。

暂把 $\cos\theta+i\sin\theta$ 简记为 $e^{i\theta}$；第~\ref{chap:ca-elementary} 章会从复指数的定义解释这个记号。由实三角函数的加法公式，
\[
 r_1e^{i\theta_1}r_2e^{i\theta_2}
 =r_1r_2e^{i(\theta_1+\theta_2)}.
\]
所以，乘以一个非零复数，相当于先按其模伸缩，再按其辐角旋转。例如乘以 $i$ 使平面逆时针旋转 $\pi/2$，乘以 $-1$ 使平面旋转 $\pi$。除法则使模相除、辐角相减；辐角的等式应按模 $2\pi$ 理解，不能直接把主值当成普通实数相加。

\begin{proposition}[棣莫弗公式与复数的根]
若 $z=re^{i\theta}\ne0$，$n$ 为正整数，则
\[
 z^n=r^ne^{in\theta},\qquad
 w^n=z\quad\Longleftrightarrow\quad
 w=r^{1/n}e^{i(\theta+2k\pi)/n},\quad k=0,1,\ldots,n-1.
\]
这 $n$ 个根两两不同，位于以原点为中心、半径为 $r^{1/n}$ 的圆周上，相邻辐角相差 $2\pi/n$。
\end{proposition}
\begin{proof}
第一式由乘法的极形式反复使用得到。若 $w=\rho e^{i\varphi}$，等式 $w^n=z$ 等价于 $\rho^n=r$ 以及 $n\varphi-\theta\in2\pi\mathbb Z$。将整数按模 $n$ 分类便得到所列的根；两个不同的 $k\in\{0,\ldots,n-1\}$ 不会给出相差 $2\pi$ 的辐角。
\end{proof}

\begin{example}[四次根与根的对称性]
求 $w^4=-16$ 的全部解。
\end{example}
\begin{solve}
写成 $-16=16e^{i\pi}$，得
\[
 w=2e^{i(\pi/4+k\pi/2)},\quad k=0,1,2,3.
\]
对应的直角坐标形式是 $\sqrt2(1+i)$、$\sqrt2(-1+i)$、$\sqrt2(-1-i)$、$\sqrt2(1-i)$。它们构成正方形。这里“求全部根”与“选定一个根函数”是不同的问题：前者给出有限集合，后者还要考虑随 $z$ 变化时能否连续选择。
\end{solve}

\begin{example}[单位根之和]
设 $n\ge2$，$\omega=\cos(2\pi/n)+i\sin(2\pi/n)$。全部 $n$ 次单位根为 $1,\omega,\ldots,\omega^{n-1}$，它们位于单位圆上等分圆周。由于 $\omega\ne1$，
\[
 1+\omega+\cdots+\omega^{n-1}=\frac{1-\omega^n}{1-\omega}=0.
\]
几何上，这说明正 $n$ 边形各顶点的位置向量之和为零。
\end{example}

\section{复平面的拓扑基础}

模给出欧氏距离 $d(z,w)=|z-w|$，因此 $\mathbb{C}$ 的距离与 $\mathbb{R}^2$ 完全相同。以下概念均相对于这一距离。

\begin{definition}[开集、闭集与边界]
若每个 $a\in U$ 都有某个圆盘 $B(a,r)\subset U$，则称 $U$ 为\emph{开集（open set）}；补集为开集的集合称为\emph{闭集（closed set）}。若 $a$ 的每个圆盘都与集合 $E$ 相交，则 $a$ 属于 $E$ 的闭包 $\overline E$；若每个圆盘同时与 $E$ 及其补集相交，则 $a$ 属于边界 $\partial E$。若每个去心圆盘都与 $E$ 相交，则称 $a$ 为 $E$ 的\emph{聚点（accumulation point）}。
\end{definition}

闭包允许用集合中的点本身检验，聚点则要求任意小邻域内还有别的点。例如单点集是闭集，但没有聚点。数列 $z_n=x_n+iy_n$ 收敛到 $z=x+iy$，当且仅当 $x_n\to x$ 且 $y_n\to y$，因为
\[
 |x_n-x|,\ |y_n-y|\le|z_n-z|\le|x_n-x|+|y_n-y|.
\]
同样的估计说明复数列是柯西列当且仅当两个坐标列都是柯西列。由实数完备性，$\mathbb{C}$ 也是完备的。

集合 $K$ 称为\emph{紧集（compact set）}，若它的每个开覆盖都有有限子覆盖。实欧氏空间中的海涅--博雷尔定理（Heine--Borel theorem）直接给出：复平面中的集合紧，当且仅当它有界且闭。例如闭圆盘紧，开圆盘不紧。若 $K\ne\varnothing$ 紧、$U$ 开且 $K\subset U$，则存在 $\delta>0$，使每个 $z\in K$ 都有 $B(z,\delta)\subset U$。证明是为每个点选取一个圆盘，先将半径减半，再对这些半径减半的圆盘取有限子覆盖。

\begin{definition}[连通与区域]
集合 $E$ 称为\emph{连通（connected）}，若不能写成两个不交、非空、在 $E$ 中相对开集合的并；这里“在 $E$ 中相对开”指形如 $E\cap U$，其中 $U$ 在 $\mathbb{C}$ 中开。非空连通开集称为\emph{区域（domain）}。若 $E$ 中任意两点可由留在 $E$ 内的连续曲线连接，则称 $E$ 道路连通（path connected）。
\end{definition}

\begin{proposition}
\label{prop:ca-polygonal-connected}
复平面中每个区域内的任意两点均可用有限折线连接。
\end{proposition}
\begin{proof}
固定 $a\in\Omega$，令 $A$ 为可由 $\Omega$ 内有限折线与 $a$ 相连的点集。对任意 $z\in\Omega$，取 $B(z,r)\subset\Omega$。若 $z\in A$，把一小段直线接在已有折线之后，可得 $B(z,r)\subset A$；若 $z\notin A$，该圆盘也不能含有 $A$ 中的点，否则可再接线段到达 $z$。故 $A$ 及 $\Omega\setminus A$ 在 $\Omega$ 中都开。由连通性及 $a\in A$，只能有 $A=\Omega$。
\end{proof}

道路连通集合必连通，因为区间的连续像连通。结合命题~\ref{prop:ca-polygonal-connected}，开集的连通性与道路连通性等价；这不是任意集合都具有的性质。

\section{曲线、同伦与单连通性}

一条\emph{曲线（curve）}是连续映射 $\gamma:[0,1]\to\mathbb{C}$。若 $\gamma(0)=\gamma(1)$，称其闭合；闭曲线若除了这对端点外不重复取值，称为简单闭曲线（simple closed curve）。积分中还要求曲线分段 $C^1$，即把参数区间有限分段后，各段上 $\gamma$ 连续可微。曲线的像是点集，曲线本身还包含行进方向和重复经过的次数，二者不能混淆。

\begin{definition}[同伦与单连通]
同端点曲线 $\gamma_0,\gamma_1:[0,1]\to\Omega$ 的\emph{定端点同伦（homotopy relative to endpoints）}是连续映射 $H:[0,1]^2\to\Omega$，满足
\[
 H(s,0)=\gamma_0(s),\quad H(s,1)=\gamma_1(s),\quad
 H(0,t)=\gamma_0(0),\quad H(1,t)=\gamma_0(1).
\]
区域 $\Omega$ 称为\emph{单连通（simply connected）}，若每条闭曲线都能在 $\Omega$ 内连续收缩为常值曲线，即存在连续 $H$，使 $H(s,0)=\gamma(s)$、$H(s,1)=a$ 且 $H(0,t)=H(1,t)$。这里中间各曲线均闭合，但公共端点可以随 $t$ 移动。
\end{definition}

若区域关于 $a$ 星形，即对每个 $z\in\Omega$，线段 $[a,z]\subset\Omega$，则 $H(s,t)=(1-t)\gamma(s)+ta$ 给出所需收缩。因此星形区域单连通，特别地，圆盘、半平面和整个复平面均单连通。圆环 $r<|z|<R$（$0<r<R$）与去心平面则不单连通：绕孔一周的曲线不能在区域内缩成一点。其严格证明将由绕数的同伦不变性给出，见第~\ref{chap:ca-integral} 章和第~\ref{chap:ca-cauchy} 章。

\begin{example}[几个典型区域]
圆盘、圆环和裂平面 $\mathbb{C}\setminus(-\infty,0]$ 都开且连通。圆环不单连通；裂平面关于点 $1$ 星形，故单连通。两个不交开圆盘的并是开集，却不连通，因此不是区域。“有界”“连通”和“单连通”是不同的要求。
\end{example}

\section{扩充复平面与球面表示}

为统一研究有限点与无穷远处的行为，在 $\mathbb{C}$ 中添加一个点 $\infty$，得到\emph{扩充复平面（extended complex plane）}$\widehat{\mathbb{C}}$。无穷远点的基本邻域定义为 $\{z:|z|>R\}\cup\{\infty\}$，其中 $R>0$；于是 $z_n\to\infty$ 当且仅当 $|z_n|\to\infty$，而不要求辐角趋向某一方向。

令 $S^2=\{(X,Y,Z)\in\mathbb{R}^3:X^2+Y^2+Z^2=1\}$，北极为 $N=(0,0,1)$。球极投影（stereographic projection）将 $S^2\setminus\{N\}$ 映到复平面：
\begin{equation}
 z=\frac{X+iY}{1-Z},\qquad
 (X,Y,Z)=\frac{1}{1+|z|^2}
 \bigl(2\operatorname{Re}z,2\operatorname{Im}z,|z|^2-1\bigr).
 \label{eq:ca-stereographic}
\end{equation}
这两个公式互为逆映射；当 $|z|\to\infty$ 时，右侧趋于 $N$。把 $N$ 对应于 $\infty$，就得到\emph{黎曼球面（Riemann sphere）}的表示。它也是紧空间：含有 $\infty$ 的一个开集覆盖某个大圆盘的外部，剩余闭圆盘可用有限个开集覆盖。

无穷远点附近的局部坐标取为 $w=1/z$，并规定 $w(\infty)=0$。因而研究 $f$ 在无穷远处的行为，可以转为研究 $f(1/w)$ 在 $w=0$ 附近的行为。$\infty$ 是新增的几何点，不是一个可任意参与域运算的复数；例如 $\infty-\infty$ 没有由此定义。

\begin{example}[阿波罗尼斯圆]
设 $a,b\in\mathbb{C}$、$a\ne b$，$k>0$。求满足 $|z-a|=k|z-b|$ 的点集。
\end{example}
\begin{solve}
平方并展开，得到
\[
 (1-k^2)|z|^2-2\operatorname{Re}\bigl(z(\overline a-k^2\overline b)\bigr)
 +|a|^2-k^2|b|^2=0.
\]
若 $k=1$，这表示到 $a,b$ 等距的垂直平分线。若 $k\ne1$，配方得圆
\[
 \left|z-\frac{a-k^2b}{1-k^2}\right|=\frac{k|a-b|}{|1-k^2|}.
\]
因此距离之比固定的轨迹可以是圆，也可以在特殊参数下退化为直线。
\end{solve}

\section{习题}
\label{sec:ca-plane-ex}
\begin{enumerate}
 \item 计算 $(1+2i)/(2-i)$，并求其模；说明结果为什么与模的商一致。
 \item 求 $z^3=8i$ 的全部根，写出极形式，并解释根之间的角度关系。
 \item 将 $|z-1|=2|z+1|$ 化为实坐标方程，指出其圆心与半径。
 \item 设 $|z|=|w|=1$ 且 $z+w\ne0$，证明 $(z-w)/(z+w)$ 为纯虚数。
 \item 判断 $\{\operatorname{Re}z>0\}$、$\{0<|z|<1\}$、$\mathbb C\setminus(-\infty,0]$ 是否为单连通区域，并说明理由。

  \item\label{ex:ca-roots-sum} 求 $z^4=-4$ 的全部解。设 $\omega$ 为本章例题中的 $n$ 次单位根，证明对任意整数 $m$，$\sum_{k=0}^{n-1}\omega^{mk}$ 等于 $n$ 或 $0$，并给出两种情形的条件。
  \item\label{ex:ca-parallelogram} 证明 $|z+w|^2+|z-w|^2=2(|z|^2+|w|^2)$，并解释其几何意义。
  \item\label{ex:ca-domains} 分别判断 $\{z:|z|<1\}$、$\{z:0<|z|<1\}$、$\mathbb{C}\setminus\mathbb{R}$ 和 $\{z:\operatorname{Re}z>0\}$ 的连通性与单连通性。
  \item\label{ex:ca-separation} 设非空紧集 $K$ 包含于开集 $U\ne\mathbb{C}$，证明 $\inf\{|z-w|:z\in K,\ w\notin U\}>0$。
\end{enumerate}

